% Problem formulation section for tech report

\subsection{Exact-Quota Width Assignment}

Let $\mathbf{h} = [h_1, \ldots, h_N]$ be a fixed real-valued FIR
coefficient vector and let $\mathcal{W} = \{16, 8, 4\}$ be the set of
allowed binary widths. Given integer quotas $n_{16}$, $n_8$, $n_4$ with
$n_{16}+n_8+n_4=N$, the problem is to find the width assignment
$\mathbf{w} \in \mathcal{W}^N$ minimizing~\eqref{eq:sse} subject to
\begin{equation}
|\{i : w_i = w\}| = n_w, \quad w \in \mathcal{W}.
\label{eq:quota}
\end{equation}
The quantization operator $Q(h,w)$ rounds $h$ to the nearest value
representable with $w$ bits (1 sign, $w-1$ fractional bits, no saturation
within the tested range).

\subsection{Equivalence to Unit-Capacity Transportation}

Problem~\eqref{eq:sse}--\eqref{eq:quota} is a binary assignment MILP.
Let $x_{i,k} \in \{0,1\}$ indicate that coefficient $i$ receives width
$\text{choices}(k)$; the cost of assignment $x_{i,k}=1$ is
$c_{i,k} = (h_i - Q(h_i, \text{choices}(k)))^2$. Then:
\begin{align}
\min_{\mathbf{x}} \quad & \sum_{i=1}^N \sum_{k=1}^3 c_{i,k} x_{i,k}
\label{eq:milp_obj}\\
\text{s.t.} \quad
& \sum_k x_{i,k} = 1, \quad i=1,\ldots,N \label{eq:milp_assign}\\
& \sum_i x_{i,k} = n_k, \quad k=1,2,3 \label{eq:milp_quota}\\
& x_{i,k} \in \{0,1\} \label{eq:milp_binary}
\end{align}
is a unit-capacity bipartite transportation / min-cost flow problem. The
network has a source connected to each of the $N$ coefficient nodes
(capacity 1, cost 0), each coefficient node connected to each of the 3
width nodes (capacity 1, cost $c_{i,k}$), and each width node connected
to a sink (capacity $n_k$, cost 0). Sending $N$ units of flow minimizes
total cost and satisfies all quotas exactly.

The LP relaxation obtained by replacing $x_{i,k}\in\{0,1\}$ with
$x_{i,k}\geq 0$ has integral extreme points: after multiplying one
partition of constraint rows by $-1$, its constraint matrix is a directed
bipartite node--arc incidence matrix and is totally unimodular. Since the
assignment and quota right-hand sides are integers, an optimal LP extreme
point is integral. Thus the assignment formulation can also be solved as a
linear program; the reported MILP runs are implementation-level
cross-checks, not evidence that integrality requires branch-and-bound.

\subsection{Reduction to Scalar Exact Integer Cost}

Llorens et al.~\cite{llorens2005integercost} study the minimization of
an L-alpha coefficient-vector error under a scalar exact integer total
cost $\sum_i r_{w_i} = R^*$, solved by a constrained-shortest-path DP.
The present three-width exact-quota problem is a special case. Set
$B = N+1$, $r_{16} = B+1$, $r_8 = 2$, $r_4 = 1$. Then for any
assignment, $R = N + B n_{16} + n_8$; because $0 \le n_8 < B$, the
value $R^* = N + B n_{16} + n_8$ uniquely encodes the quotas
$(n_{16}, n_8, n_4)$. The Llorens DP at exact cost $R^*$ therefore
solves the quota-constrained problem. This was empirically verified on
all 108 saved test configurations (maximum objective discrepancy:
floating-point rounding).

The DP developed in Section~\ref{sec:algorithm} is a direct state
encoding that tracks used quota counts rather than scalar cost; it is
equivalent in result but distinct in state representation. Neither
formulation constitutes a new optimization class.


\noindent\textbf{Scope of the quotas.}\quad
The counts constrain coefficient representations, not measured LUT,
DSP, power, or throughput budgets. The selected 16/8/4-bit quotas are
experimental settings; this study does not establish a product-level
requirement for those exact counts or a fixed inventory of physical
multipliers. Hardware synthesis may remove zero constants or map equal
counts differently. Solving the allocation problem therefore does not
solve the broader hardware resource--response design problem.
